\documentclass{article}

\usepackage[utf8]{inputenc}     % pdfLaTeX
\usepackage[T1]{fontenc}
\usepackage[magyar]{babel}

\usepackage{amsthm}
\usepackage{amsmath}
\usepackage{amsfonts}
\usepackage{tikz}
\usetikzlibrary{arrows.meta, positioning}
\usepackage{graphicx}

\title{Kombinatorikai feladatok megoldása kódolással}
\author{Rovenszky Robin}
\date{Legutóbb frissítve: 2026. Március 23.}

\begin{document}

\maketitle
\tableofcontents
\newpage

\section{Első feladat}
Hányféleképpen írhatjuk fel a tizet három nem negatív egész szám összegeként, ha az összeadandók sorrendje is számít?

\section{Második feladat}
A huszat hányféle képpen lehet felírni három pozitív egész szám összegeként, ha az összeadandók sorrendje is számít?

Általánosan $n \in \mathbb{N^+}$ számot $k \in \mathbb{N^+}$ darab számból 
\[
n-1 \choose k-1
\]
féleképpen írhatunk fel, $1 \le k \le n$.\\\\
Ha $k \in \mathbb{N}$, akkor pedig
\[
n + k - 1 \choose k - 1
\]
féleképpen.

\end{document}